\section{引言}
\label{sec:intro}

线上部署的大语言模型智能体通过周期性强化学习更新持续进化：每隔一段时间，用新采集的交互数据精炼策略。这本质上是\textbf{持续强化学习}问题——模型必须在习得新能力的同时保住已有能力。与实验室一次性集齐全分布数据不同，线上数据更新快、到达不均衡，往往只能\textbf{按能力域分批、顺序地}训练。这恰好放大了\textbf{灾难性遗忘}~\cite{D1}：后训练的域会系统性侵蚀先训练的域已掌握的能力。

针对遗忘的主流解法是\textbf{经验回放}~\cite{A2}：训练新数据时从历史经验采样一部分回放进当前批次，用独立的持续学习损失让模型温习旧能力。已有工作的隐含假设是被回放的经验携带忠实的学习信号。在多能力域场景下，一个额外挑战浮现：不同能力会在扁平缓冲区中互相挤占，"回放什么、回放多少"本身就是一个设计问题。

本文面向智能体大语言模型的持续强化学习，设计并实现一套完整的训练与评测系统，核心贡献有三：

\begin{enumerate}
  \item \textbf{按能力分区的九桶经验回放缓冲区}。桶按能力域而非难度划分，每桶设固定容量上限（任务数平方根加权）和保护下限。桶内淘汰、禁止跨桶挤出，桶间按能力空间\textbf{距离}采样回放——离当前训练批越远的桶遗忘风险越高、回放权重越大。桶内采样和淘汰策略为可替换设计，当前分别采用随机采样和FIFO淘汰。

  \item \textbf{面向持续学习的复合训练目标}，以零框架改动注入GRPO：策略梯度、对参考策略的逆向KL锚定、经验回放损失、固定开启的熵正则（防策略多样性坍缩~\cite{B4}）。奖励由单个外部冻结大模型裁判在统一细则下打分，多维乘性聚合，安全作为二元门控以提高评价一致性。

  \item \textbf{按桶分组的抗遗忘评测协议}：全部能力域顺序训完后一次统一评测，施加按训练先后衰减的权重后置量化"越早训练的桶遗忘越重"，并与同规模无持续学习机制的基线对照。
\end{enumerate}

本文在Qwen3.5-9B智能体上跨九个能力域实现并验证整套系统，用Pass$^3$基准~\cite{C1}同时刻画新能力习得与旧能力遗忘。
